\chapter{Correlation-kernel density criterion}\label{ch:correlation-kernel}

\section{A second analytic route from the Xi kernel}

The canonical Fourier-kernel representation introduced in the operator programme admits a second route that is independent of the Hermite--Biehler candidate audited in XF-2.  For the even Riemann kernel $\Phi$ and
\[
\Xi(z)=\int_{\mathbb R}\Phi(t)e^{izt}\,\dd t,
\]
define the second correlation function
\begin{equation}\label{eq:nu2-correlation}
\nu_2(t)
=
\int_{-\infty}^{\infty}
(t-2s)^2\Phi(t-s)\Phi(s)\,\dd s.
\end{equation}
For $0<|y|<1/2$, set
\begin{equation}\label{eq:phi2y-correlation}
\Phi_{2,y}(t)=\cosh(ty)\,\nu_2(t).
\end{equation}
These are the $n=2$ correlation objects used by Dimitrov and Xu in their Fourier--Wronskian criteria \cite{DimitrovXu2016}.

\section{Published RH-equivalent density criterion}

Dimitrov--Xu Theorem~1.1 supplies a standard external equivalence.  In the present notation,
\begin{equation}\label{eq:dx-density-rh}
\boxed{
\mathrm{RH}
\iff
\forall\,y\in\left(-\frac12,\frac12\right)\setminus\{0\},\quad
\overline{\operatorname{span}\mathcal T(\Phi_{2,y})}^{\,L^1(\mathbb R)}
=L^1(\mathbb R),
}
\end{equation}
where $\mathcal T(f)$ denotes the family of all translates of $f$.

The same work supplies an equivalent convolution-annihilator formulation: for every admissible $y$, the condition
\begin{equation}\label{eq:dx-annihilator}
\Phi_{2,y}*g=0,\qquad g\in L^\infty(\mathbb R),
\end{equation}
forces $g=0$ exactly on the RH surface \cite{DimitrovXu2016}.

Within the typed implication graph, the theorem itself has status \status{STANDARD}.  The global density condition and the equivalent annihilator condition have status \status{OPEN}; they are the mathematical premises that must be independently established before the standard equivalence can close the RH node.

\section{Wronskian bridge}

For $n=2$, the Fourier--Wronskian identity gives
\begin{equation}\label{eq:dx-wronskian}
W_2(\Xi;x)
=
\Xi(x)\Xi''(x)-\Xi'(x)^2
=
-\mathcal F[\nu_2](x),
\end{equation}
with the Fourier convention of \cite{DimitrovXu2016}.  Since $\Xi$ is real on the real $z$ axis, its first Laguerre quantity obeys
\begin{equation}\label{eq:laguerre-wronskian}
\boxed{
L_1[\Xi](x)
=
\Xi'(x)^2-\Xi(x)\Xi''(x)
=
-W_2(\Xi;x)
=
\mathcal F[\nu_2](x).
}
\end{equation}
This gives an exact consistency bridge between the correlation-kernel and derivative representations.

\section{Executable diagnostic layer}

The implementation in \texttt{src/critical\_axis/correlation\_kernel.py} evaluates $\nu_2$, $\Phi_{2,y}$, the Wronskian, and the Laguerre quantity with explicit finite truncation controls.  The validation suite checks:
\begin{enumerate}[label=(\roman*)]
\item evenness and positive reference values of the implemented kernel;
\item numerical evenness of $\nu_2$;
\item the exact $y\leftrightarrow-y$ symmetry of $\Phi_{2,y}$;
\item real-axis closure of Eq.~\eqref{eq:laguerre-wronskian};
\item fail-closed enforcement of $0<|y|<1/2$;
\item sampled Laguerre positivity as a \status{NUMERICAL\_DIAGNOSTIC} quantity.
\end{enumerate}

Finite samples and finite translate families carry diagnostic authority only.  The repository records this layer as
\[
\status{NUMERICAL\_DIAGNOSTIC}.
\]
The theorem-level target retains its full quantifiers over every admissible $y$, all translates, and the complete space $L^1(\mathbb R)$.

\section{Relation to the branch-cancellation programme}

The two active routes now share the same canonical Xi kernel but expose different closure obligations:
\begin{align}
\Phi
&\longrightarrow A_\pm
\longrightarrow \text{exact zero cancellation}
\longrightarrow \text{population/coordinate bridge},\\
\Phi
&\longrightarrow \nu_2
\longrightarrow \Phi_{2,y}
\longrightarrow \text{$L^1$ translation density or annihilator closure}.
\end{align}
The first route isolates a local zero-state representation problem.  The second route converts RH into a global harmonic-analysis completeness problem.  Their coexistence supplies a useful cross-audit: a future mechanism that closes the critical axis should be checked against both the local branch structure and the global correlation-kernel criterion.

\section{Current frontier}

The immediate XF-3 proof obligation is therefore
\begin{equation}\label{eq:xf3-frontier}
\boxed{
\forall\,0<|y|<\frac12:\quad
\overline{\operatorname{span}\mathcal T(\Phi_{2,y})}^{\,L^1(\mathbb R)}
=L^1(\mathbb R),
}
\end{equation}
or equivalently the global bounded-annihilator condition of Eq.~\eqref{eq:dx-annihilator}.  The solver records both as \status{OPEN\_RH\_EQUIVALENT\_CRITERION}; the Riemann-hypothesis node therefore retains status \status{OPEN}.
